\section{Khảo sát hiện trạng}

\subsection{Nhu cầu học tiếng Anh của trẻ em hiện nay}

Trong bối cảnh công nghệ ngày càng phát triển và quá trình hội nhập quốc tế diễn ra mạnh mẽ, tiếng Anh đã trở thành một kỹ năng quan trọng đối với trẻ em từ sớm. Không chỉ được sử dụng trong học tập, tiếng Anh còn xuất hiện phổ biến trong nhiều lĩnh vực như công nghệ, truyền thông, giải trí và giao tiếp toàn cầu. Vì vậy, nhiều phụ huynh hiện nay mong muốn cho con tiếp cận ngoại ngữ từ nhỏ nhằm nâng cao khả năng giao tiếp và thích nghi với môi trường hiện đại.


Theo các nghiên cứu về giáo dục, giai đoạn từ 6 đến 10 tuổi được xem là thời điểm trẻ có khả năng tiếp thu ngôn ngữ hiệu quả. Trong độ tuổi này, trẻ dễ ghi nhớ âm thanh, bắt chước cách phát âm và hình thành phản xạ ngôn ngữ nhanh hơn. Nếu được học tập bằng phương pháp phù hợp, trẻ có thể tiếp cận tiếng Anh một cách tự nhiên và hứng thú hơn.


Tuy nhiên, việc học tiếng Anh của nhiều trẻ hiện nay vẫn còn tồn tại một số hạn chế. Phần lớn trẻ chủ yếu học theo hình thức truyền thống như ghi nhớ từ vựng hoặc làm bài tập ngữ pháp, khiến quá trình học trở nên khô khan và dễ gây nhàm chán. Nhiều em tuy có thể làm bài tập tốt nhưng khả năng nghe, nói và phát âm vẫn chưa thực sự hiệu quả.


Ngoài ra, sự phổ biến của điện thoại thông minh và internet cũng đã làm thay đổi thói quen học tập của trẻ nhỏ. Thay vì chỉ tiếp cận kiến thức qua sách vở, trẻ hiện nay thường hứng thú hơn với hình ảnh, video, trò chơi và các ứng dụng mang tính tương tác trên thiết bị di động. Điều này tạo điều kiện thuận lợi cho việc phát triển các ứng dụng học tiếng Anh trực quan, sinh động và phù hợp với trẻ em.


Xuất phát từ những lý do trên, việc xây dựng một ứng dụng hỗ trợ học tiếng Anh cho trẻ em kết hợp công nghệ AI nhằm tăng khả năng tương tác và hỗ trợ luyện phát âm là hướng phát triển phù hợp với nhu cầu học tập hiện nay.


\subsection{Thực trạng học tiếng Anh truyền thống của trẻ em}

Hiện nay tại nhiều trường và trung tâm ngoại ngữ, học tiếng Anh theo kiểu truyền thống mang tính lý thuyết và thiên về ghi nhớ. Mặc dù đã có nhiều cải tiếng trong giáo dục nhưng việc học tiếng ở nhiều nơi vẫn còn nhiều hạn chế nhất định cần phải cải được cải thiện để phù hợp xu hướng trẻ em ngày nay.


\noindent\textbf{a) Phương pháp học còn nặng về lý thuyết}


Nhiều chương trình học hiện nay vẫn còn tập trung quá nhiều vào ngữ pháp, các cấu trúc câu và bài tập viết. Trẻ thường phải học thuộc từ vựng và làm các bài kiểm tra trên giấy mà lại có ít có cơ hội thực hành giao tiếp thực tế.


Đối với các bạn trẻ nhỏ, việc học theo cách ghi nhớ một cách máy móc rất dễ gây nhàm chán. Trẻ thường mất tập trung nhanh nếu bài học không có hình ảnh minh họa, âm thanh hoặc các hoạt động tương tác.


Ngoài ra, việc học theo kiểu khuôn mẫu cũng khiến trẻ khó phát triển phản xạ ngôn ngữ tự nhiên. Mặc dù có nhiều bạn trẻ có thể làm bài tập tốt nhưng lại không tự tin khi nghe hoặc nói tiếng Anh.


\noindent\textbf{b) Thiếu môi trường luyện phát âm}


Phát âm là một trong những kỹ năng quan trọng nhất khi học ngoại ngữ. Tuy nhiên, trong môi trường học truyền thống, giáo viên khó có thể sửa phát âm cho từng học sinh do số lượng học sinh trong lớp đông, mình cô không thể hướng dẫn cho từng bạn phát âm chuẩn được.


Có nhiều trẻ phát âm sai trong thời gian dài nhưng lại không có ai giúp chỉnh sửa kịp thời, dẫn đến việc hình thành thói quen phát âm không chuẩn. Điều đó sẽ ảnh hưởng lớn đến kỹ năng nghe và giao tiếp của trẻ về sau này.


Ngoài ra,các bạn trẻ thường rất ngại nói tiếng Anh trước đám đông vì sợ phát âm sai hoặc bị bạn chê cười. Điều này làm giảm khả năng phản xạ ngôn ngữ của trẻ, thiếu đi sự tư tin để học và cải thiện phát âm.


\noindent\textbf{c) Thiếu tính trự quan và sinh động}


Trẻ em sẽ học hiệu quả hơn khi được tiếp xúc với hình ảnh, âm thanh và các trò chơi tương tác. Tuy nhiên, nhiều phương pháp học truyền thống vẫn sử dụng giáo trình khô khan, ít hình ảnh minh họa hoặc thiếu các hoạt động tương tác.


Các bài học được kéo dài liên tục sẽ dễ khiến trẻ mất hứng thú. Trong khi đó, trẻ nhỏ lại có xu hướng tiếp thu tốt hơn thông qua hoạt động vui chơi và trải nghiệm.


\noindent\textbf{d) Hạn chế về thời gian và địa điểm học}


Ở các lớp học truyền thống thường cố định về thời gian và địa điểm trong một thời gian dài. Điều này gây khó khăn cho phụ huynh trong việc đưa đón con em đi học, cũng dễ khiến các bạn trẻ rời vào vòng lặp chán nản quanh quẩn việc đến trường.


Ngoài ra, trẻ em chỉ được học trong khoảng thời gian nhất định nên khó duy trì việc luyện tập hàng ngày. Trong khi đó, việc học ngoại ngữ cần được thực hành thường xuyên để đạt hiệu quả tốt.


\subsection{Hiện trạng các ứng dụng học tiếng Anh trên thị trường}

Trong những năm gần đây, với sự phát triển manhh mẽ của công nghệ số thị trường ứng dụng học tiếng Anh phát triển rất mạnh. Có rất nhiều ứng dụng hỗ trợ học từ vựng, luyện nghe, luyện nói và học giao tiếp dành cho nhiều đối tượng khác nhau.


Một số ứng dụng nổi bật hiện nay gồm:


\begin{itemize}
\item Monkey Junior
\item Duolingo
\item Lingokids
\item Elsa Speak
\item Cake
\item BBC Learning English
\end{itemize}

Các ứng dụng trên đã góp phần thay đổi cách tiếp cận việc học ngoại ngữ của nhiều người dùng. Tuy nhiên, sau khi khảo sát thực tế, có thể nhận thấy rằng nhiều ứng dụng vẫn chưa thực sự phù hợp với trẻ em Việt Nam trong độ tuổi nhỏ.


\begin{figure}[H]
\centering
\includegraphics[width=0.85\textwidth]{app_monkey_junior.jpg}
\caption{Hình ảnh app Monkey Junior}
\end{figure}

\noindent\textbf{a) Nội dung chưa phù hợp hoàn toàn với trẻ em}


Một số ứng dụng được thiết kế cho nhiều độ tuổi vì vậy nội dung chưa tối ưu riêng cho trẻ từ 6 đến 10 tuổi. Nhiều bài học có cấu trúc quá phức tạp hoặc sử dụng từ vựng quá khó khiến trẻ tiếp cận khó và lại dễ gây thêm chán nản cho trẻ.


Bên cạnh lại có những ứng dụng sử dụng giao diện quá nhiều chữ, ít hình ảnh chưa phù hợp hoặc màu sắc chưa đủ sinh động để thu hút trẻ em.


\noindent\textbf{b) Nhiều tính năng bị giới hạn}


Hiện nay, có nhiều ứng dụng yêu cầu người dùng phải trả phí để sử dụng đầy đủ các tính năng như luyện nói AI, mở khóa bài học hoặc để theo dõi tiến độ học tập. Điều này gây khó khăn đối với nhiều phụ huynh khi muốn cho con học lâu dài.


\noindent\textbf{c) Chưa tối ưu khả năng luyện phát âm tiếng Anh}


Một số ứng dụng có hỗ trợ nhận diện giọng nói nhưng chưa đánh giá chính xác phát âm của trẻ em Việt Nam. Kết quả chấm điểm đôi khi chưa sát với thực tế hoặc phản hồi còn chậm.


Ngoài ra, nhiều ứng dụng chỉ tập trung vào luyện từ đơn lẻ mà chưa hỗ trợ phản xạ giao tiếp một cách tự nhiên.


\noindent\textbf{d) Thiếu sự kết hợp giữa học và giải trí}


Trẻ em thường thích học thông qua trò chơi và hình ảnh sinh động. Tuy nhiên, nhiều ứng dụng hiện nay vẫn còn thiên về học lý thuyết nên chưa tạo được cảm giác hứng thú lâu dài cho trẻ.


\begin{figure}[H]
\centering

\includegraphics[width=0.85\textwidth]{app_duolingo.jpg}
\caption{Hình ảnh app duolingo}
\end{figure}

\subsection{Khảo sát nhu cầu thực tế của phụ huynh và trẻ em}

Trong quá trình tìm hiểu tình hình các ứng dụng thực tế, nhu cầu của phụ huynh đối với một ứng dụng học tiếng Anh cho trẻ em tập trung vào các yếu tố sau:


\begin{itemize}
\item Đối với phụ huynh
\item Giao diện đơn giản, dễ sử dụng
\item Nội dung phù hợp với trẻ nhỏ
\item Có hỗ trợ luyện phát âm chuẩn
\item Có thể học mọi lúc mọi nơi
\item Theo dõi được kết quả học tập của con
\item Có yếu tố giải trí để trẻ hứng thú hơn
\item Đối với trẻ em
\item Thích hình ảnh sinh động
\item Thích học qua trò chơi
\item Thích âm thanh và hoạt hình
\item Muốn tương tác giống trò chuyện thực tế
\item Thích được chấm điểm và nhận phần thưởng
\end{itemize}

Từ kết quả khảo sát thực tế cho thấy, một ứng dụng học tiếng Anh hiệu quả cần kết hợp giữa giáo dục và giải trí, đồng thời tận dụng các công nghệ AI để tăng khả năng tương tác với người học.


\section{Yêu cầu hệ thống:}

\subsection{Mục tiêu xây dựng hệ thống}

Hệ thống ``Tiếng Anh Cho Bé'' được xây dựng nhằm hỗ trợ trẻ em học tiếng Anh theo phương pháp trực quan, sinh động và hiện đại.


Mục tiêu chính của hệ thống gồm:


\begin{itemize}
\item Giúp trẻ học tiếng Anh hiệu quả hơn thông qua hình ảnh và trò chơi
\item Hỗ trợ luyện phát âm bằng AI
\item Tăng khả năng phản xạ tiếng Anh cho trẻ
\item Tạo môi trường học tập vui vẻ và dễ tiếp cận
\item Xây dựng hệ thống học tập có thể mở rộng trong tương lai
\end{itemize}

\subsection{Đối tướng sử dụng hệ thống}

Hệ thống được hướng đến các đối tượng sử dụng chính như sau:


\noindent\textbf{a) Trẻ em từ 6 đến 10 tuổi}


Đây là đối tượng sử dụng chính của ứng dụng Android và Web User. Các bạn trẻ sẽ sử dụng hệ thống để:


\begin{itemize}[leftmargin=1.3cm,itemsep=0.3cm]
	\item Học từ vựng.
	\item Nghe phát âm.
	\item Luyện nói.
	\item Chơi Quiz.
	\item Chat AI.
	\item Dịch ảnh tiếng Anh.
\end{itemize}


\noindent\textbf{b) Phụ huynh}


Phụ huynh sẽ hỗ trợ trẻ trong quá trình tạo tài khoản, học tập và theo dõi kết quả học của con.


\noindent\textbf{c) Quản trị viên (Admin)}


Admin sẽ sử dụng hệ thống quản trị Web để:


\begin{itemize}[leftmargin=1.3cm,itemsep=0.3cm]
	\item Quản lý chủ đề học.
	\item Quản lý từ vựng.
	\item Quản lý câu hỏi quiz.
	\item Quản lý kết quả quiz
	\item Quản lý tài khoản người dùng.
	\item Theo dõi kết quả học tập.
\end{itemize}



\subsection{Phân tích chức năng hệ thống}

\begin{itemize}[label=$\diamond$,leftmargin=1.2cm,itemsep=0.35cm]
	
	\item \textbf{Chức năng đăng ký và đăng nhập}
	
	Cho phép người dùng tạo tài khoản và truy cập hệ thống. Hệ thống sẽ kiểm tra dữ liệu hợp lệ trước khi lưu vào cơ sở dữ liệu và truy cập vào hệ thống.
	\item \textbf{Chức năng học từ vựng}
	
	Người dùng sẽ chọn các chủ đề học như màu sắc, động vật, trái cây,...
	
	Mỗi từ vựng sẽ bao gồm:
	
	\begin{itemize}[label=\textbullet,leftmargin=1.2cm,itemsep=0.15cm]
		\item Hình ảnh minh họa
		\item Phiên âm
		\item Nghĩa tiếng Việt.
	\end{itemize}
	
	\item \textbf{Chức năng luyện phát âm AI}
	
	Người dùng đọc từ tiếng Anh như mẫu, hệ thống sẽ dùng AI Speech Recognition để:
	
	\begin{itemize}[label=\textbullet,leftmargin=1.2cm,itemsep=0.15cm]
		\item Nhận diện giọng nói
		\item So sánh với phát âm chuẩn
		\item Đánh giá độ chính xác
		\item Chấm điểm phát âm
	\end{itemize}
	
	\item \textbf{Chức năng chat AI}
	
	Cho phép người dùng đặt câu hỏi bằng tiếng Anh hoặc tiếng Việt, AI sẽ phản hồi tự động cho người dùng:
	
	\begin{itemize}[label=\textbullet,leftmargin=1.2cm,itemsep=0.15cm]
		\item Luyện giao tiếp
		\item Học từ mới
		\item Tăng phản xạ ngôn ngữ
	\end{itemize}
	
	\item \textbf{Chức năng chụp ảnh dịch}
	
	Người dùng chụp ảnh chứa nội dung là tiếng Anh, hệ thống sẽ dịch nghĩa sang tiếng Việt tự động cho người dùng.
	
	\begin{itemize}[label=\textbullet,leftmargin=1.2cm,itemsep=0.15cm]
		\item Nhận diện văn bản bằng OCR
		\item Dịch sang tiếng Việt bằng AI
		\item Hiển thị kết quả trực tiếp cho người dùng
	\end{itemize}
	
	\item \textbf{Chức năng quiz game}
	
	Giúp cho trẻ ôn tập kiến thức thông qua các trò chơi bằng câu hỏi trắc nghiệm.
	
	\item \textbf{Đối với quản trị}
	
	Admin quản lý toàn bộ dữ liệu hệ thống thông qua web client.
	
\end{itemize}

\subsection{Yêu cầu chức năng}

Hệ thống được xây dựng gồm ba nhóm chức năng chính như sau:

\begin{itemize}[label=$\diamond$,leftmargin=1.2cm,itemsep=0.35cm]
	
	\item \textbf{Đối với App Android}
	
	\begin{itemize}[label=\textbullet,leftmargin=1.2cm,itemsep=0.15cm]
		\item Đăng ký tài khoản.
		\item Đăng nhập.
		\item Học từ vựng theo chủ đề.
		\item Luyện phát âm.
		\item Chat AI.
		\item Chụp ảnh dịch.
		\item Làm Quiz.
		\item Xem thành tích.
		\item Quản lý tài khoản.
		\item Đăng xuất.
	\end{itemize}
	
	\item \textbf{Đối với Web Admin (CRUD)}
	
	\begin{itemize}[label=\textbullet,leftmargin=1.2cm,itemsep=0.15cm]
		\item Đăng nhập.
		\item Quản lý người dùng.
		\item Quản lý chủ đề học.
		\item Quản lý từ vựng.
		\item Quản lý câu hỏi Quiz.
		\item Quản lý kết quả Quiz.
		\item Quản lý kết quả học tập.
		\item Quản lý lịch sử luyện nói.
		\item Đăng xuất.
	\end{itemize}
	
	\item \textbf{Đối với Web User}
	
	\begin{itemize}[label=\textbullet,leftmargin=1.2cm,itemsep=0.15cm]
		\item Đăng ký tài khoản.
		\item Đăng nhập vào hệ thống.
		\item Chọn chủ đề học.
		\item Làm Game Quiz.
		\item Luyện nói.
		\item Quản lý tài khoản.
		\item Đăng xuất.
	\end{itemize}
	
\end{itemize}

\subsection{Yêu cầu phi chức năng}

\noindent\textbf{a) Giao diện}


Thân thiện với trẻ em


Màu sắc sinh động


Hình ảnh rõ ràng


Dễ thao tác


\noindent\textbf{b) Hiệu năng}


Phản hồi nhanh


Tải dữ liệu ổn định


Hoạt động ứng dụng tốt trên Android


\noindent\textbf{c) Bảo mật}


Bảo vệ tài khoản người dùng


Mã hóa dữ liệu đăng nhập


Phân quyền Admin/User rõ ràng


\noindent\textbf{d) Khả năng mở rộng}


Hệ thống có thể mở rộng thêm:


Bài học mới


Tính năng AI nâng cao


Học giao tiếp
